\chapter{个性化候选召回}
\label{chap:rs-retrieval}

已知用户最近浏览过登山包、水壶和相机，并不意味着可以逐一比较目录中的每件商品。若目录有数百万件物品，而后续模型只能细看几百件，系统必须先决定哪些对象有机会进入比较。本章研究这一候选召回问题：把请求时可用的需求与偏好依据转化为一个规模受限、对象有效且具有足够覆盖的候选集合。

前章已经得到用户、对象和匹配的表示。本章继续追问表示如何真正取回物品：邻域表返回相邻对象，向量索引返回高分对象，树模型沿目录路径取得叶子，标识生成器则预测可以查回物品的代码。不同入口形成候选后，还要处理重复、失效与遗漏。训练分数是否适合目录竞争、评价是否让模型看到了同样的对象，是贯穿这些方法的两个检查点。

% Outline: 3.1
\section{召回任务与方法分类}
\label{sec:rs-retrieval-task}

设全目录为$I$，请求时点$t$可用的对象为$I_t\subseteq I$，最终候选为$C_t\subseteq I_t$，大小不超过$K$。展示列表$L_t$还需从候选中选择和组织，因此进入$C_t$只是取得继续比较的机会。对象资格、检索命中和最终价值是三件事：一款缺货背包可能与用户非常匹配，但没有当前购买资格；一款有效背包可能因检索遗漏而未进入候选；召回命中也不能保证排序后会展示，更不能保证用户满意。

候选取得有两个常用接口。目录检索沿已有索引、邻域或评分结构取出物品；标识生成预测一段物品代码，再通过映射恢复目录对象。这里按最后怎样取得物品分类。例如，扩散模型生成一个连续向量后查近邻，仍使用目录检索接口；语言模型输出物品代码再查表，才使用标识生成接口。两者可以组合。

匹配依据则是另一个坐标。交互与协同证据利用用户之间、物品之间的行为关系；内容和外部关系证据利用文本、图像、属性及知识；多源方法结合这些证据。点击标签只说明预测目标来自点击，不能单凭它断定输入用了什么协同关系。物品图若由共同购买构造，边仍是交互证据；若由品牌和部件关系构造，边才来自相应的属性或知识。

\begin{table}[htbp]
  \centering\small
  \begin{tabularx}{\linewidth}{@{}lXXX@{}}
    \toprule
    候选接口 & 交互与协同 & 内容与外部关系 & 多源结合 \\
    \midrule
    目录检索 & 邻域、潜在因子 & 内容近邻、关系访问 & 联合表示和多路检索 \\
    标识生成 & 历史到物品码 & 内容代码、文本条件 & 历史与语义代码结合 \\
    \bottomrule
  \end{tabularx}
  \caption{取得对象的接口与匹配证据可以交叉。同一实现可以占据多列，也可以组合两种接口；训练使用生成目标不决定候选属于哪一行。}
  \label{tab:rs-retrieval-axes}
\end{table}

本章把$K$称为候选名额，并将时延、运算量与存储作为计算预算。名额不是广告消费金额，也不是每次必须填满的指标。如果只有80件对象满足约束，返回80件有效候选比用20件无效对象凑满100件更符合任务。另一方面，只返回分数最集中的一类物品可能漏掉另一种需求；因此需要同时检查覆盖与代价。

% Outline: 3.2
\section{候选生成方法}
\label{sec:rs-candidate-methods}

一个候选模型必须同时交代输入、学习目标、分数和访问方式。只给出用户向量还不够：物品偏置能否被索引支持，多个兴趣如何分配名额，生成的代码是否对应真实对象，都会改变实际候选。以下先从能够直接访问目录的实现展开，再比较标识生成。

% Outline: 3.2.1
\subsection{基于目录检索的召回}
\label{sec:rs-directory-retrieval}

目录入口可以很简单。请求中的价格上限先过滤属性，内容表示找描述接近的物品，经过核验的关系可以找兼容部件或同主题内容。词项匹配的BM25见第~\ref{chap:rs-search}~章；近邻索引的构建与服务属于系统卷。本节关心由什么证据产生什么分数，以及如何用该分数取得目录候选。

下列模型的机制并不互斥。多个兴趣向量可以共享独立的物品编码器，图传播得到的向量可以交给近邻索引，交互集合重建也可以直接给全目录打分。前章的跨域、群组及地点模型若已输出合法的用户或群组—物品分数，同样可以在其规定的目录中取Top-$K$；迁移它们时须保留原来的成员、时间和空间条件。

% Outline: 3.2.1.1
\subsubsection{UserCF：用户邻域召回}
\label{sec:rs-usercf}

一位用户尚未评价某件物品时，可以参考评价习惯相近的用户。\term{基于用户的协同过滤}（user-based collaborative filtering, UserCF）先找到这样的邻居，再汇总邻居对未见物品的评价。GroupLens的经典实现提供了用户评分聚合的依据，但其完整新闻系统并不等同于今日的召回服务\scite{rs-h07}

显式评分中，有的人习惯打高分，有的人打分严格。令$r_{u,i}$为已观测评分，$\bar r_u$为用户已观测评分均值，$O_{u,v}$为两人共同评价的物品。可用中心化余弦形式的Pearson相关系数衡量共同变化：
\begin{equation}
 w_{u,v}=
 \frac{\sum_{i\in O_{u,v}}(r_{u,i}-\bar r_u)(r_{v,i}-\bar r_v)}
 {\sqrt{\sum_{i\in O_{u,v}}(r_{u,i}-\bar r_u)^2}
  \sqrt{\sum_{i\in O_{u,v}}(r_{v,i}-\bar r_v)^2}}.
 \label{eq:rs-usercf-sim}
\end{equation}
这里的均值取各自全部已观测评分；也可在共同集合上重新中心化，但须一致执行。没有共同观测或分母为零时，相似度没有定义，应排除该对或采用已说明的回退，不能当成满分相似。有效邻居集合$N_u$截断为固定数量。若只保留正相似邻居，预测为
\begin{equation}
 \widehat r_{u,i}=\bar r_u+
 \frac{\sum_{v\in N_u:\,r_{v,i}\text{已观测}}w_{u,v}(r_{v,i}-\bar r_v)}
 {\sum_{v\in N_u:\,r_{v,i}\text{已观测}}|w_{u,v}|}.
 \label{eq:rs-usercf-predict}
\end{equation}
只有评价过$i$的邻居进入该物品的分子和分母；分母为零时没有邻域预测。

\begin{table}[htbp]
 \centering\small
 \begin{tabular}{@{}lrrrr@{}}
 \toprule
 用户 & 登山包$a$ & 水壶$b$ & 相机$c$ & 通勤包$d$\\
 \midrule
 $u$ & 5&3&---&---\\
 $v$ & 5&3&5&1\\
 $w$ & 4&2&1&5\\
 $z$ & 1&5&2&---\\
 \bottomrule
 \end{tabular}
 \caption{用于邻域计算的教学评分表。“---”表示未观测，不是零分。下一节只把有评分视为发生交互，以展示显式与隐式口径的区别。}
 \label{tab:rs-cf-example}
\end{table}

在表~\ref{tab:rs-cf-example}中，$\bar r_u=4,\bar r_v=3.5,\bar r_w=3$。由共同的$a,b$可得$w_{u,v}=2/\sqrt5\approx0.894$，$w_{u,w}=1$，而$w_{u,z}<0$。保留$v,w$后，相机得分约为$4+(0.894\times1.5-2)/1.894=3.652$，通勤包得分约为$4+(-0.894\times2.5+2)/1.894=3.875$。召回访问邻居的已评分物品并集，滤掉用户已见的$a,b$，于是$d$排在$c$前。

这个例子也暴露了小样本问题：仅两件共同物品就可能产生相关系数1。可以要求最小共现数，或以$|O_{u,v}|/(|O_{u,v}|+\lambda)$收缩相似度，但应把它写成额外的可靠性处理。隐式行为下，改用二值或加权交互向量及余弦等度量，不能再把一次点击当作显式高评分。热门物品让许多人形成共同观测，也可能主导邻居和候选；更新用户邻域的代价还会随活跃用户与交互变化增长。

% Outline: 3.2.1.2
\subsubsection{ItemCF：物品邻域召回}
\label{sec:rs-itemcf}

若用户数很大，而单次历史较短，可以提前为每件物品保存相邻物品。\term{基于物品的协同过滤}（item-based collaborative filtering, ItemCF）从共同用户的评价或交互计算物品相似度，请求时沿历史物品的邻域访问候选\scite{rs-h03}

显式评分可采用调整余弦：对共同评价物品$i,j$的用户集合$U_{i,j}$，以$r_{u,i}-\bar r_u$与$r_{u,j}-\bar r_u$计算余弦，消除用户评分尺度。若只有隐式交互，令$U_i$为与$i$交互的用户集合，一种基础选择为
\begin{equation}
 w_{i,j}=\frac{|U_i\cap U_j|}{\sqrt{|U_i|\,|U_j|}},
 \qquad s(u,j)=\sum_{i\in H_u}a_{u,i}w_{i,j}.
 \label{eq:rs-itemcf}
\end{equation}
其中$a_{u,i}$是历史权重，可以为1，也可由已定义的次数或时间衰减给出。每个$i$只存最相近的$m$件物品时，请求最多产生约$|H_u|m$条访问结果，再按物品汇总、过滤与截断。显式评分预测常另做加权归一化；式~\eqref{eq:rs-itemcf}是候选累积分数，不是评分量表上的预测值。

仍用表~\ref{tab:rs-cf-example}，但本次只将“有评分”视为一次交互。$|U_a|=|U_b|=4$，$|U_c|=3$，$|U_d|=2$，所以$w_{a,c}=w_{b,c}=3/\sqrt{12}\approx0.866$，$w_{a,d}=w_{b,d}=2/\sqrt8\approx0.707$。用户$u$的两条历史均取权重1，则相机累计1.732，通勤包累计1.414。这里$c$领先，原因是丢弃评分高低后，共现证据发生了变化，不能把它理解成UserCF与ItemCF必然持有相反偏好。

物品邻域表可以离线预计算，共现计数或评分变化后再更新。它减少了请求时的用户比较，却没有免除邻域维护。来自交互的相近也不表示内容相同：相机与存储卡可能是互补关系。新物品没有共同用户时得不到可靠邻域，热门物品可能在大量历史路径中反复出现，需要通过覆盖与增量候选诊断。

% Outline: 3.2.1.3
\subsubsection{MF：潜在因子检索}
\label{sec:rs-mf}

邻域表依赖可观察的重叠；\term{矩阵分解}（matrix factorization, MF）则让用户和物品共享一个低维空间，用少量潜在因子解释评分。这个空间联系了没有直接共现的用户和物品，但每个坐标未必对应可命名的属性\scite{rs-h08}

设$\bm{h}_u,\bm{v}_i\in\mathbb{R}^d$，全局均值为$\mu$，用户和物品偏置为$b_u,b_i$。只在已观测集合$\Omega$上拟合
\begin{align}
 \widehat r_{u,i}&=\mu+b_u+b_i+\bm{h}_u^\top\bm{v}_i,\\
 \mathcal L&=\frac12\sum_{(u,i)\in\Omega}
  (r_{u,i}-\widehat r_{u,i})^2+
  \frac\lambda2\left(\sum_u(\|\bm h_u\|^2+b_u^2)
  +\sum_i(\|\bm v_i\|^2+b_i^2)\right).
 \label{eq:rs-mf-objective}
\end{align}
缺失项不进入平方误差，不能填零后当成评分0。对一条样本，记$e=r_{u,i}-\widehat r_{u,i}$；在局部采用正则化随机更新时，因子按
$\bm h_u\leftarrow\bm h_u+\eta(e\bm v_i-\lambda\bm h_u)$、
$\bm v_i\leftarrow\bm v_i+\eta(e\bm h_u-\lambda\bm v_i)$更新，右侧应使用更新前的两个向量。按样本分摊全局正则时还需相应设置每次的正则系数；偏置也由误差和正则项更新。

教学例取$\mu=3,b_u=0,b_i=0.1$，$\bm h_u=(0.6,0.2)$，$\bm v_i=(0.4,0.5)$。预测为3.44；若观测评分4，则$e=0.56$。取$\eta=0.1,\lambda=0.01$作一次局部更新，得到$\bm h'_u=(0.6218,0.2278)$，$\bm v'_i=(0.4332,0.5107)$。误差通过同一用户或同一物品的共享参数，影响随后其他配对的预测。

用于目录召回时，$\mu+b_u$对同次请求是常数，但$b_i$不是。要严格按该预测取Top-$K$，可增广$\widetilde{\bm h}_u=(\bm h_u,1)$和$\widetilde{\bm v}_i=(\bm v_i,b_i)$，以增广内积访问支持该目标的索引；也可直接扫描目录。先按纯内积截断再加偏置，只是近似，可能漏掉高偏置对象。MF因此既能预测评分，也能比较给定候选，或在完整目录上召回；三种使用方式不能仅凭“训练了因子”就视为已经实现。成对偏好训练的BPR由第~\ref{chap:rs-ranking}~章展开。

% Outline: 3.2.1.4
\subsubsection{隐式反馈MF：置信度加权}
\label{sec:rs-implicit-mf}

观看次数没有统一的评分刻度。\term{隐式反馈MF}把“是否发生行为”与“有多少证据”分开：令原始次数为$n_{u,i}$，偏好目标$p_{u,i}=\mathbb I[n_{u,i}>0]$，置信度$c_{u,i}=1+\alpha n_{u,i}$。次数增长增加拟合权重，而不是把偏好目标从1线性抬高\scite{rs-h09}

其目标在全部用户—物品对上定义：
\begin{equation}
 \mathcal L=\sum_{u,i}c_{u,i}(p_{u,i}-\bm h_u^\top\bm v_i)^2
 +\lambda\left(\sum_u\|\bm h_u\|^2+\sum_i\|\bm v_i\|^2\right).
 \label{eq:rs-implicit-mf}
\end{equation}
固定物品矩阵$\bm V$（每行为$\bm v_i^\top$），令$\bm C_u$为置信度对角阵，$\bm p_u$为目标向量。对$\bm h_u$求导并令梯度为零，得到
\begin{equation}
 \bm h_u=(\bm V^\top\bm C_u\bm V+\lambda\bm I)^{-1}
             \bm V^\top\bm C_u\bm p_u.
 \label{eq:rs-als}
\end{equation}
实际以线性方程求解代替显式求逆。固定用户侧后对物品作同样更新，就是交替最小二乘（alternating least squares, ALS）。由于$\bm C_u-\bm I$仅在交互项上非零，$\bm V^\top\bm C_u\bm V$可以拆成共享的$\bm V^\top\bm V$和稀疏修正，不必为每个用户重复累加整个目录。

若观看次数为$(3,0,1)$且$\alpha=2$，则目标为$(1,0,1)$，置信度为$(7,1,3)$。作为一维计算例，固定物品因子均为1，$\lambda=1$，式~\eqref{eq:rs-als}给出用户因子$(7+3)/(7+1+3+1)=5/6$。未观测项确实以低权重把预测向0约束；这是模型目标的约定，不是用户明确拒绝。所得内积可直接用于目录候选。与显式MF相比，损失覆盖了未观测项；与BPR相比，它拟合加权点目标，而非要求正例超过负例。频繁观看、误触和被动播放仍需由数据语义区分。

% Outline: 3.2.1.5
\subsubsection{YouTube DNN：候选生成}
\label{sec:rs-youtube-candidate}

因子表对已有用户方便，却不容易直接接收新发生的历史和请求特征。\term{YouTube候选生成网络}把观看、搜索及其他请求信息汇成一个用户向量，再与视频输出嵌入匹配。原模型是用户网络配合物品嵌入表，不是两端都采用深层网络的对称双塔\scite{rs-r01}

观看和搜索词嵌入经过汇总，与人口统计等特征拼接，通过多层网络得到$\bm h(\bm x)$。以下一观看物品$i^+$作分类标签，理想目录目标为
\begin{equation}
 p_\theta(i\mid\bm x)=
 \frac{\exp(\bm h(\bm x)^\top\bm v_i)}
 {\sum_{j\in I}\exp(\bm h(\bm x)^\top\bm v_j)}.
 \label{eq:rs-directory-softmax}
\end{equation}
大目录训练使用采样近似，具体竞争集合和采样修正见第~\ref{sec:rs-sampled-softmax}~节。训练要按预测位置组织历史，防止把目标后的信息放进输入。服务计算一次用户向量，按物品内积近邻取回候选，因而不需要为每个物品重跑用户网络。

假设历史含登山与摄影视频，汇总输入经教学网络得到$(0.8,0.3)$。下一观看视频为“背包打包”，其向量$(1,0)$；两个训练对照的向量为$(0,1)$和$(0.2,0.2)$，得分为0.8、0.3和0.22。这三个分数形成一轮训练竞争，部署时却要面对完整视频目录，而非永远只比较这两项负例。论文的独立排序网络还会接收更丰富的候选特征，其训练目标和本候选网络应分开理解。

% Outline: 3.2.1.6
\subsubsection{DSSM：独立编码与语义匹配}
\label{sec:rs-dssm}

只有物品ID时，新描述不能自然进入因子空间。\term{深层结构语义模型}（deep structured semantic model, DSSM）让两端各自编码文本，再比较它们的语义表示。原始任务是搜索查询与点击文档匹配，采用字符三元组散列压缩输入词汇，经网络变换得到向量，以余弦相似度的条件似然学习\scite{rs-h10}

令两端表示为$\bm h=f_\theta(q)$、$\bm v=g_\phi(d)$，相似度$c(q,d)=\bm h^\top\bm v/(\|\bm h\|\|\bm v\|)$。以正点击文档$d^+$和采样对照构成竞争集合$D$，最小化
\begin{equation}
 \ell_{\mathrm{DSSM}}=-\log
       \frac{\exp(\gamma c(q,d^+))}
            {\sum_{d\in D}\exp(\gamma c(q,d))}.
 \label{eq:rs-dssm-loss}
\end{equation}
其中$\gamma$控制相似度在softmax中的尺度。文档编码不依赖查询，因此可提前保存；归一化后的内积可直接用于向量检索。

将这一接口迁移到推荐时，必须重定义两端。一个教学改造以“通勤，防水，预算500元”的需求文本编码用户条件，以商品资料编码物品，并用实际交互或人工匹配标签训练。假设两端单位向量分别为$(1,0)$与$(0.8,0.6)$，余弦为0.8；另一商品$(0,1)$得0，前者可进入语义候选。但价格等硬条件仍应独立核验，文本相近也不等于该用户更愿意购买。这个例子说明独立编码的可复用性，不是原DSSM论文报告的推荐实验。

% Outline: 3.2.1.7
\Needspace{90mm}
\subsubsection{MIND：多兴趣召回}
\label{sec:rs-mind}

一个用户同时关注运动和烹饪时，单向量可能落在两类物品之间。\term{MIND}以多个兴趣槽分别汇总行为，让每个槽独立取回候选。动态路由指行为与槽之间的分配随当前匹配反复更新，而非提前用类别标签划分历史\scite{rs-r29}

设变换后的行为向量为$\bm z_j$，路由logit为$b_{j,k}$。一次迭代计算行为$j$分给槽$k$的权重$c_{j,k}=\operatorname{softmax}_k(b_{j,k})$，再得到
\begin{equation}
 \bm a_k=\sum_j c_{j,k}\bm z_j,\qquad
 \bm h_k=\frac{\|\bm a_k\|^2}{1+\|\bm a_k\|^2}
                  \frac{\bm a_k}{\|\bm a_k\|},\qquad
 b_{j,k}\leftarrow b_{j,k}+\bm z_j^\top\bm h_k.
 \label{eq:rs-mind-routing}
\end{equation}
零向量的压缩结果定义为零。原方法采用行为共享的映射及随机路由初始化，避免所有槽始终保持对称。教学例中，运动行为为$(1,0)$，烹饪行为为$(0,1)$，初始两行分配分别为$(0.8,0.2)$、$(0.2,0.8)$。第一槽汇总$(0.8,0.2)$，压缩为约$(0.393,0.098)$；运动行为向第一槽的logit增量0.393，大于向第二槽的0.098，于是后续迭代可强化分工。

路由输出还与用户资料嵌入拼接，经共享的多层变换得到最终兴趣向量$\widetilde{\bm h}_k$。训练时已知正物品，标签感知注意以目标嵌入计算各兴趣权重，用目标相关的汇总向量进入采样softmax。原文通过对匹配值作幂变换调节注意集中程度，并实验采用近似硬选择；不能把MIND一概说成始终平均多个兴趣。目标参与的是训练中的兴趣选择，不是线上已知事实。

服务去掉标签感知注意，由每个$\widetilde{\bm h}_k$分别访问目录近邻，再合并。例如两槽各取3件，最多形成6个不同物品，也可能因重复只有4件。原方法还按历史长度启发式地限制有效槽数，让短历史使用较少向量。槽数是模型容量和计算选择，不是测量出的真实兴趣数；是否覆盖两类需求要看最终候选。

% Outline: 3.2.1.8
\subsubsection{ComiRec：多兴趣与可控聚合}
\label{sec:rs-comirec}

多兴趣接口确定后，还可以分别研究兴趣怎样提取、候选怎样聚合。\term{ComiRec-DR}采用动态路由，ComiRec-SA则用自注意力生成多组历史权重。后一种形式可写为
$\bm A=\operatorname{softmax}_{\mathrm{history}}(\bm W_2\tanh(\bm W_1\bm H^\top))$，
再以$\bm A\bm H$得到多行兴趣向量；每行权重在历史位置上归一化\scite{rs-r05}

训练选与正物品内积最大的兴趣向量计算推荐目标，可与MIND调节标签感知注意的做法比较；MIND也可趋近硬选择，不能仅按“硬／软”将两者完全对立。服务同样让每个兴趣取回近邻。固定运动、烹饪历史和总名额6时，可让两兴趣各取3件，再观察是否某一路返回大量重复；不能一边增加向量和名额，一边把候选改善全部归因于兴趣提取。

可控聚合在候选上兼顾相关性与类别差异。例如设候选$a,b,c$相关分数为0.9、0.85、0.8，$a,b$同类而$c$异类；选择$a$后，若异类增益权重为0.1，下一步$b$的增益为0.85，$c$为0.9，就会选$c$。这是教学化的边际增益计算，说明类别项怎样改变集合；它不保证异类一定对应用户的另一种需求。候选聚合见第~\ref{sec:rs-candidate-fusion}~节，完整列表的效用由第~\ref{chap:rs-presentation}~章建立。

% Outline: 3.2.1.9
\subsubsection{PinSage：内容与图关系召回}
\label{sec:rs-pinsage}

物品内容与用户整理行为可以互补。\term{PinSage}在物品—收藏板二部关系上通过随机游走选取重要邻居，再结合内容学习可检索的物品表示；收藏关系来自用户整理行为，不能改称外部知识边\scite{rs-r03}

对物品$i$，游走访问频率形成归一化邻居权重$\alpha_{i,j}$。一层聚合先变换邻居特征并加权池化，得到$\bm n_i=\sum_j\alpha_{i,j}f(\bm v_j)$，再将自身表示与$\bm n_i$拼接、变换并归一化。训练用相关正物品与负物品约束相对距离或内积，原方法采用间隔式排序并结合困难负例，使表示适合相关物品检索。

设两个收藏板分别包含“登山包—水壶”和“登山包—相机”。假设游走对水壶、相机的权重为0.75和0.25，变换后特征为$(1,0)$、$(0,1)$，则邻居消息为$(0.75,0.25)$。它还需经过自身融合才成为最终表示，不可把这一步直接当成完成的图网络。服务可预计算物品向量，以种子物品取近邻；用于个人推荐时另规定用户历史中的种子及汇总权重。训练使用图并不意味着请求时在线遍历全图，内容相近、共收藏和个体偏好也仍需区分。

% Outline: 3.2.1.10
\subsubsection{LightGCN：协同图传播召回}
\label{sec:rs-lightgcn}

纯交互场景并不一定需要每层复杂的特征变换。\term{LightGCN}在用户—物品二部图上传播初始ID嵌入，以对称度归一化控制高连接节点的贡献\scite{rs-r30}
\begin{equation}
 \bm h_u^{(\ell+1)}
 =\sum_{i\in N_u}\frac{\bm v_i^{(\ell)}}{\sqrt{|N_u||N_i|}},
 \qquad
 \bm v_i^{(\ell+1)}
 =\sum_{u\in N_i}\frac{\bm h_u^{(\ell)}}{\sqrt{|N_i||N_u|}}.
 \label{eq:rs-lightgcn}
\end{equation}
最终表示是包含第0层的多层加权和，评分为最终用户与物品的内积。传播没有特征变换矩阵和逐层非线性，但多跳邻居仍然把协同关系引入表示。原模型用BPR训练初始嵌入，最终向量可预计算后用于目录近邻。

用表~\ref{tab:rs-cf-example}的交互图，用户$u$连接$a,b$，两个物品各有4位用户。若第0层物品向量为$\bm v_a=(1,0),\bm v_b=(0,1)$，则$\bm h_u^{(1)}=(1,1)/\sqrt8\approx(0.354,0.354)$。若自身第0层为$(1,0)$，等权组合两层得到$(0.677,0.177)$。完整候选分数还须使用同样完成传播和组合的物品向量；只替换用户侧并不是完整LightGCN。

与MF相比，这些因子受到图传播的结构耦合；与PinSage相比，节点类型、邻居选取、内容输入和训练目标都不同。原LightGCN仅有ID和交互，零度新物品没有可传播的证据，不能靠“使用图网络”就宣称解决冷启动。传播深度还可能使表示趋同，并增加更新范围。

% Outline: 3.2.1.11
\Needspace{90mm}
\subsubsection{TDM：层级目录的深层检索}
\label{sec:rs-tdm}

向量内积便于索引，但用户与候选的复杂交互不一定能分离成两端向量。\term{树深度模型}（tree-based deep model, TDM）让深层评分器比较用户与树节点，通过逐层裁剪控制访问规模\scite{rs-r59}叶子是物品，内部节点表示一组后代，训练把正物品的祖先作为各层正节点，并与同层负节点比较。

考虑八件物品的二叉目录。第一层节点$A,B$分别管前四件、后四件；第二层为$A_1,A_2,B_1,B_2$；再下一层到物品。束宽取2时，第一层保留两节点，第二层假设得分为0.9、0.6、0.8、0.3，则只展开$A_1,B_1$，其四个叶子再竞争最终名额。评分可以读取用户历史和节点特征，并不要求固定内积。

图~\ref{fig:rs-tree-retrieval}标出了这一裁剪发生的位置。被舍弃的内部节点仍对应真实物品，只是本次请求不再访问其后代。
\begin{figure}[htbp]
 \centering
 % 八物品教学目录；节点得分仅用于TDM束宽为2的演示。
\begin{tikzpicture}[
 font=\figurefont\footnotesize,
 branch/.style={rectangle,draw=BookRule,fill=BookPaper,
   align=center,minimum width=12mm,minimum height=8mm,inner sep=3pt},
 kept/.style={branch,draw=FigureModel,fill=FigureModel!10},
 leaf/.style={branch,minimum width=8mm,minimum height=7mm},
 edge/.style={draw=BookMuted,line width=.8pt},
 active/.style={draw=FigureModel,line width=1.1pt},
 cut/.style={draw=BookMuted,dashed,line width=.8pt}
]
 \node[branch] (root) at (0,0) {目录};
 \node[kept] (a) at (-2.6,-1.1) {$A$};
 \node[kept] (b) at (2.6,-1.1) {$B$};
 \node[kept] (a1) at (-3.9,-2.3) {$A_1$\\0.9};
 \node[branch] (a2) at (-1.3,-2.3) {$A_2$\\0.6};
 \node[kept] (b1) at (1.3,-2.3) {$B_1$\\0.8};
 \node[branch] (b2) at (3.9,-2.3) {$B_2$\\0.3};
 \node[leaf,fill=FigureOutput!10] (i1) at (-4.55,-3.6) {$a$};
 \node[leaf,fill=FigureOutput!10] (i2) at (-3.25,-3.6) {$b$};
 \node[leaf] (i3) at (-1.95,-3.6) {$c$};
 \node[leaf] (i4) at (-.65,-3.6) {$d$};
 \node[leaf,fill=FigureOutput!10] (i5) at (.65,-3.6) {$e$};
 \node[leaf,fill=FigureOutput!10] (i6) at (1.95,-3.6) {$f$};
 \node[leaf] (i7) at (3.25,-3.6) {$g$};
 \node[leaf] (i8) at (4.55,-3.6) {$h$};
 \draw[active] (root) -- (a);
 \draw[active] (root) -- (b);
 \draw[active] (a) -- (a1);
 \draw[edge] (a) -- (a2);
 \draw[active] (b) -- (b1);
 \draw[edge] (b) -- (b2);
 \draw[active] (a1) -- (i1);
 \draw[active] (a1) -- (i2);
 \draw[active] (b1) -- (i5);
 \draw[active] (b1) -- (i6);
 \draw[cut] (a2) -- (i3);
 \draw[cut] (a2) -- (i4);
 \draw[cut] (b2) -- (i7);
 \draw[cut] (b2) -- (i8);
 \node[align=center,text=BookMuted] at (0,-4.4)
   {第二层保留$A_1,B_1$；虚线路径不再展开};
\end{tikzpicture}

 \caption{八物品教学目录上的逐层裁剪。第二层按假设分数保留两个节点，只继续比较$a,b,e,f$；虚线表示本次未展开的路径。图示目录也可用于理解SEATER的路径标识，但后者的生成概率不等于图中的TDM节点分数。}
 \label{fig:rs-tree-retrieval}
\end{figure}

当分支数与束宽固定、树保持平衡时，约$\log |I|$层访问可控制评分次数，但这不是精确全目录Top-$K$保证。若$A_2$内部有一个得分很高的叶子，它仍可能因祖先被删而永远不可见。TDM以层级组织、训练目标及树学习减少这种错失，无法在任意评分函数下消除它。树、节点表示和训练样本需随目录组织同步维护，见第~\ref{sec:rs-candidate-maintenance}~节。

\Needspace{85mm}
% Outline: 3.2.1.12
\subsubsection{PinRec：结果条件下的生成向量召回}
\label{sec:rs-pinrec}

同一用户在搜索、信息流与相关物品页面留下的历史可以共同利用，但当前入口和希望支持的行为不同。\term{PinRec}以包含物品、查询、动作与时间的跨入口序列训练生成模型，再用目标入口的曝光日志微调。输出头接受入口和目标行为条件，产生与物品空间匹配的稠密向量，随后访问近邻目录\scite{rs-r71}

训练以真实后续交互所对应的物品向量作为匹配目标，用带批内采样频率修正的softmax学习内积匹配，令共享历史状态支持不同结果条件。训练中的输入使用真实历史，条件来自目标结果；服务时才从多个结果条件的输出中抽取一个向量反馈到序列，再继续预测下一步表示。服务将多个步骤、多个行为条件的向量按名额配置发给近邻检索，过于相近的向量可以合并并累加名额。生成器承担候选提供，后面的排序和混排仍然存在。

例如，假设点击条件分配3个候选名额，收藏条件分配2个。两个教学向量分别检索到$(a,b,c)$与$(b,d)$，去重后得到$\{a,b,c,d\}$，其中$b$同时具有两条来源。如果希望最终取得5件不同物品，还需要续取或名额重分配。把更多名额放在收藏条件下会改变比较机会，但不保证用户真的收藏；名额也不是广告资金预算。多步向量的收益须连同额外解码、近邻查询和重复率一起评价。

% Outline: 3.2.1.13
\subsubsection{MBGCN：多行为图的目标行为召回}
\label{sec:rs-mbgcn}

购买较少时，浏览与收藏能补充交互，但把边合并会丢失行为含义。\term{MBGCN}分别利用用户—物品传播和按行为定义的物品—物品传播。前者学习哪些行为对该用户更有帮助，后者学习同一种行为连接了哪些物品\scite{rs-r98}

设用户$u$的行为$b$出现$n_{u,b}$次，可学习的行为强度为$w_b$，则行为权重为$\alpha_{u,b}=w_bn_{u,b}/\sum_cw_cn_{u,c}$。模型先对各行为邻居作均值聚合，再按$\alpha_{u,b}$组合并线性变换；物品侧汇总相邻用户。另一条传播在“物品—同一行为用户—物品”的关系上，使用行为专属变换得到$\bm z_{i,b}$。各层表示拼接后，完整分数可记为
\begin{equation}
 s(u,i)=\rho\,\bm h_u^\top\bm v_i+
 (1-\rho)\sum_b\left(
 \frac1{|H_{u,b}|}\sum_{j\in H_{u,b}}\bm z_{j,b}\right)^\top
 \bm M_b\bm z_{i,b},
 \label{eq:rs-mbgcn-score}
\end{equation}
空行为集合的贡献定义为0。$\bm M_b$学习该行为下的物品关系，$\rho$权衡两类分数。模型用目标行为正物品与未观测物品的BPR损失训练，浏览输入不自动成为购买正标签。

假设用户浏览相机、购买背包，各一次；$w_{\text{浏览}}=1,w_{\text{购买}}=3$，则用户消息分别占0.25和0.75。若某候选的用户内积为0.6，两种历史关系项合计0.2，$\rho=0.7$，完整分数为0.48；另一候选内积0.5、关系项0.6，分数0.53，反而领先。只按第一项取回候选会改变模型目标。

服务可以提前保存图表示，并汇总用户每类历史向量。式~\eqref{eq:rs-mbgcn-score}仍可写成增广内积：把$\rho\bm h_u$与各$(1-\rho)\bm M_b^\top\bar{\bm z}_{u,b}$拼接，对应拼接$\bm v_i$和$\bm z_{i,b}$后检索；这是由完整评分式得到的访问实现，不要求每次遍历全部历史。高频浏览噪声和稀疏关系仍可能影响表示。多行为输入与多头响应预估的区别见第~\ref{chap:rs-ranking}~章。

% Outline: 3.2.1.14
\Needspace{90mm}
\subsubsection{DiffNet：社交邻域传播的候选匹配}
\label{sec:rs-diffnet}

个人记录不足时，关注或信任关系提供了另一种邻域。\term{DiffNet}把用户ID因子与属性融合为初始状态，再逐层聚合可信邻居的上一层状态，与自身状态拼接后变换。物品表示由物品ID与属性融合得到；原DiffNet的社交传播发生在用户侧，没有DiffNet++中的物品传播扩展\scite{rs-r99}

令社交邻居为$S_u$，一层计算为
$\bm h_u^{(\ell+1)}=f_\ell([\operatorname{Pool}_{v\in S_u}\bm h_v^{(\ell)};\bm h_u^{(\ell)}])$。
得到末层状态后，还加入本人已交互物品的均值：
\begin{equation}
 \widetilde{\bm h}_u=\bm h_u^{(L)}
       +\frac1{|H_u|}\sum_{j\in H_u}\bm v_j,\qquad
 s(u,i)=\widetilde{\bm h}_u^\top\bm v_i.
 \label{eq:rs-diffnet}
\end{equation}
对空历史需单独约定回退。训练用正交互和采样未观测对象构成成对排序监督；服务预计算图状态与历史均值，再以完整用户向量访问目录。

设甲关注乙，乙关注丙，三人初始一维状态为0、0、1。为了只观察消息路径，假设更新为自身与唯一邻居各占一半。一层后甲仍为0，乙为0.5；两层后甲才收到丙的信息，变为0.25。真实模型的变换可学习，但两跳信息不能在一次同步一层更新中凭空抵达。关系与历史变化还会使缓存状态过时。社交相似可能来自共同背景或主动选择朋友，预测中的传播不能识别真实社会影响；错误关注关系和过深聚合都可能损害匹配。

% Outline: 3.2.1.15
\subsubsection{MHCN：高阶用户关系的多通道匹配}
\label{sec:rs-mhcn}

两两邻接无法单独表达“互有联系的三人”与“共同购买的两人”之间的结构差别。\term{MHCN}从社交边和交互构造局部模式，分成社交、联合、购买三个通道：社交通道使用三用户的社交三角；联合通道连接有社交关系且购买同一物品的用户；购买通道保留共同购买但没有相应社交连接的关系\scite{rs-r100}

这里的购买模式含两个用户和一个物品，不应把“三人共同购买”直接画成原模型规定的三用户购买三角。若甲、乙、丙都买过背包，它产生多个用户对的共同购买关系；其中甲乙互相关注，则这一对进入联合关系，其余无社交连接的配对进入购买关系。图中的物品帮助确定模式，用户通道的传播矩阵仍作用于用户。

各通道先通过门控选择初始用户特征，再按模式诱导的归一化邻接$\bm D_c^{-1}\bm A_c$传播，并组合不同层。通道注意权重$\alpha_{u,c}$将它们融合；额外的用户—物品图传播补充物品信息，最终得到用户和物品向量，以内积及BPR学习推荐。作为聚合例，甲的两个社交通道邻居状态为$(1,0)$与$(0,1)$且权重相同，则消息为$(0.5,0.5)$。若三个通道输出依次为$(0.5,0.5),(0.8,0.2),(0.2,0.8)$，注意权重为$(0.2,0.5,0.3)$，融合结果为$(0.56,0.44)$；还需与协同传播分量组合后才能评分。

S$^2$-MHCN在推荐损失之外增加用户、用户周围子超图与全图层次的互信息目标，避免多通道融合丢失关系信息。服务依赖最终向量，图模式构建与刷新是额外成本。辅助目标没有证明邻居持有相同心理偏好；本任务仍给个人产生候选，不是前章的群组共同选择。

% Outline: 3.2.1.16
\subsubsection{Mult-VAE：交互集合重建的目录评分}
\label{sec:rs-mult-vae}

只要用户已有少量行为，就可以由共享编码器直接计算其潜在状态。\term{Mult-VAE}输入目录维度的隐式交互向量$\bm x_u$，编码器输出高斯后验的均值和方差，解码器把潜变量$\bm z_u$映射为整个目录的logits。多项分布似然要求历史交互物品获得较多概率质量\scite{rs-r105}变分自编码器的通用推导见\BookChapterRef{chap:pgm-generative}{模型卷的“概率生成模型”章}，这里重点说明其目录输出与推荐数据含义。

设$\bm\pi_\theta(\bm z)=\operatorname{softmax}(f_\theta(\bm z))$，训练最大化
\begin{equation}
 \mathbb E_{q_\phi(\bm z\mid\widetilde{\bm x}_u)}
       \left[\sum_i x_{u,i}\log\pi_{\theta,i}(\bm z)\right]
 -\beta D_{\mathrm{KL}}\!\left(q_\phi(\bm z\mid\widetilde{\bm x}_u)
                              \,\|\,\mathcal N(\bm0,\bm I)\right).
 \label{eq:rs-mult-vae}
\end{equation}
输入丢弃产生$\widetilde{\bm x}_u$，重建目标仍是原交互；训练可逐步增加并选择KL权重$\beta$。服务关闭输入扰动，取潜在均值解码，再按目录logits过滤与截断。概率在全目录上归一化，表达相对分配，不是每件物品独立的点击率。

教学输入为$(1,0,1,0)$，输出logits假设为$(2,1,0,3)$。softmax约为$(0.237,0.087,0.032,0.644)$；过滤已见的第一、第三项后，第四项领先于第二项。这里既没有物品码解码，也不必预先保存每位用户的自由参数，但输出层仍有目录维度，计算和物品增删都受其影响。集合输入不保留顺序，空历史缺乏个体化信息，只能依赖共享先验或额外特征。

% Outline: 3.2.1.17
\subsubsection{RecVAE：复合先验与交替训练的目录评分}
\label{sec:rs-recvae}

稀疏输入下，共享编码器的一次更新可能改善一批用户，却使另一批用户的后验大幅变化。\term{RecVAE}在Mult-VAE接口上加入旧编码器约束，让当前表示既受总体先验限制，也保留此前已学的信息\scite{rs-r106}其复合先验的基本形式是
\begin{equation}
 p(\bm z\mid\bm x_u)=
 \alpha\mathcal N(\bm z;\bm0,\bm I)
 +(1-\alpha)q_{\phi_{\mathrm{old}}}(\bm z\mid\bm x_u).
 \label{eq:rs-recvae-prior}
\end{equation}
旧参数在当前优化阶段固定。原文实验还向混合分布加入较宽的高斯分量；混合的是概率密度，不能把均值加权后当作同一个高斯，再套用两高斯KL的闭式公式。

例如在某个采样点，标准高斯密度为0.2，旧后验密度为0.5，$\alpha=0.2$，则基本复合先验的密度为0.44。当前后验在该点的对数密度与$\log0.44$之差进入KL的采样估计。模型还令KL权重与用户交互数量成比例，编码器在扰动输入下学习，解码器更新关闭输入丢弃；两组参数交替优化，旧编码器按训练轮次刷新。

若输入$(1,0,1,0)$被扰动为$(1,0,0,0)$，编码器仍需支持原来两项交互的重建，旧后验则从未扰动输入计算。这并不是每次线上请求重新训练：服务仍只读取当前编码器，得到潜在状态并解码目录分数。复合先验解决的是训练稳定与正则选择，不自动修正曝光偏差；固定目录维度、空历史和顺序丢失的限制仍在。

% Outline: 3.2.1.18
\subsubsection{DiffRec：交互分布去噪的目录评分}
\label{sec:rs-diffrec}

交互集合也可通过逐步加噪和恢复学习。\term{DiffRec}以目录维度的历史向量$\bm x_0$为起点，控制噪声强度以保留个体信息；去噪网络同时读取带噪向量和扩散步数，预测原交互\scite{rs-r107}扩散模型的通用概率结构见\BookChapterRef{chap:pgm-generative}{模型卷的“概率生成模型”章}。用$\ell$区分扩散步数与请求时点，前向状态为
\begin{equation}
 \bm x_\ell=\sqrt{\bar\alpha_\ell}\bm x_0+
       \sqrt{1-\bar\alpha_\ell}\bm\epsilon,\qquad
 \bm\epsilon\sim\mathcal N(\bm0,\bm I).
 \label{eq:rs-diffrec-noise}
\end{equation}
其中$\alpha_\ell=1-\beta_\ell$为该步保留比例，$\bar\alpha_\ell=\prod_{j=1}^\ell\alpha_j$为累计系数。训练的第一步使用原向量平方重建误差；后续步的误差权重为
$\tfrac12[\bar\alpha_{\ell-1}/(1-\bar\alpha_{\ell-1})-\bar\alpha_\ell/(1-\bar\alpha_\ell)]$。
该权重来自反向高斯分布的匹配，不能把各步一概当成相同监督强度。

教学例取$\bm x_0=(1,0)$、$\bar\alpha_2=0.81$、一次噪声取值$\bm\epsilon=(0,1)$，得到$(0.9,\sqrt{0.19})\approx(0.9,0.436)$。反向均值使用
\begin{equation}
 \bm\mu_\theta(\bm x_\ell,\ell)=
 \frac{\sqrt{\alpha_\ell}(1-\bar\alpha_{\ell-1})}{1-\bar\alpha_\ell}\bm x_\ell
 +\frac{\sqrt{\bar\alpha_{\ell-1}}(1-\alpha_\ell)}{1-\bar\alpha_\ell}
       \widehat{\bm x}_0.
 \label{eq:rs-diffrec-mean}
\end{equation}
若$\alpha_2=\bar\alpha_1=0.9$，网络预测$\widehat{\bm x}_0=(0.8,0.2)$，两个系数均约0.499，恢复结果约为$(0.849,0.318)$。这组假设只展示计算，不说明训练能恢复到什么精度。

服务从已有历史的带噪状态开始。原方法允许先加噪$T'$步，将所得状态置于反向链的第$T$步，再执行$T$次确定性均值更新；通常$T'\le T$，因此不能把前向加噪步数直接当作反向调用次数。所得目录分数再过滤已见和不合格对象。L-DiffRec通过物品聚类与潜在压缩降低维度，T-DiffRec在交互权重中纳入时间。压缩会改变表示，多步恢复也增加推断成本，须与全目录输出一起衡量。去噪目标中的“噪声”是模型构造的扰动，并不证明日志中的误触、曝光偏差或偏好变化都被消除了。

\Needspace{90mm}
% Outline: 3.2.1.19
\subsubsection{DreamRec：历史条件下的连续目标向量生成}
\label{sec:rs-dreamrec}

重建整个目录向量成本较高；\term{DreamRec}改为恢复下一目标物品的连续嵌入。Transformer编码已知历史，扩散训练向真实下一物品嵌入加噪，并让网络在历史条件下预测原嵌入。训练随机移除条件，使同一网络能够给出有条件和无条件预测，服务再以无分类器引导调节历史条件的影响\scite{rs-r108}

设历史条件为$\bm h$，有条件与无条件的目标预测分别为$f_\theta(\bm z_\ell,\ell,\bm h)$和$f_\theta(\bm z_\ell,\ell,\varnothing)$。一种等价的引导记号为
$\widehat{\bm z}_0=(1+w)f_\theta(\bm z_\ell,\ell,\bm h)-w f_\theta(\bm z_\ell,\ell,\varnothing)$。
它增强两种预测之差；$w$改变条件强度，不能保证越大越符合偏好。生成从噪声开始，经过恢复得到连续向量，再按与真实物品嵌入的内积取Top-$K$。

若教学生成结果为$(1,0.5)$，背包与相机的目录向量为$(1,0)$、$(0,1)$，内积分数为1与0.5，候选来自目录中的背包和相机，而不是把向量本身当成一件新商品。训练的目标向量重建不需要推荐负采样，但目标仍由已观察到的下一交互给出，因此没有消除观测机制的影响。多步恢复、条件强度与物品表示刷新共同影响召回；“理想物品”只是预测向量的解释，不能作为真实心理对象的证明。

% Outline: 3.2.2
\subsection{基于标识生成的召回}
\label{sec:rs-identifier-retrieval}

目录索引也可以换成可生成的寻址代码。设物品$i$的标识为$c(i)=(c_1,\ldots,c_m)$，生成器读入历史和请求条件$\bm x$，学习
\begin{equation}
 \begin{aligned}
 p_\theta(c(i)\mid\bm x)
 &=\prod_{\ell=1}^m p_\theta(c_\ell\mid\bm x,c_{<\ell}),\\
 \mathcal L_{\mathrm{ID}}
 &=-\sum_{\ell=1}^m\log p_\theta(c_\ell(i^+)\mid\bm x,c_{<\ell}(i^+)).
 \end{aligned}
 \label{eq:rs-id-generation}
\end{equation}
训练使用正目标的真实前缀，服务则依据已生成前缀续写。束搜索每一步只保留有限数量的高分前缀，最终得到若干标识，再恢复物品。词元词表中的任意组合未必对应目录对象，因而代码概率与有效物品概率不能无条件等同。

这种接口需要共同设计四件事：如何编码对象、生成器读取什么、何种监督影响代码与生成、如何把输出落回目录。表~\ref{tab:rs-id-models}给出各实现的接口差别。量化码本可以由许多物品共享，但代码—物品映射仍随目录增长；树节点与文本索引又有不同存储范围。训练与维护成本不能只统计生成网络。

\begin{table*}[htbp]
 \centering\small
 \renewcommand{\arraystretch}{1.12}
 \begin{tabularx}{\linewidth}{@{}lXXX@{}}
 \toprule
 模型 & 标识与输入依据 & 主要学习安排 & 候选恢复与维护 \\
 \midrule
 TIGER & 内容残差量化码、交互历史 & 先编码，后预测标识 & 束搜索、消歧后查表 \\
 OneRec-Think & 物品码与文字 & 文本对齐、理由监督、反馈优化 & 理由条件下生成码，核验目录 \\
 DCGR & 结果前缀与语义码 & 日志监督结果及物品 & 联合束搜索，业务偏置调节 \\
 LC-Rec & 文本量化码与历史 & 均匀映射，多类语言对齐 & 代码解码，维护唯一映射 \\
 LETTER & 内容与协同量化码 & 重构、协同、多样性约束 & 冻结代码后生成与映射 \\
 SEATER & 协同向量构成平衡树路径 & 路径预测、层次对比与排序 & 合法路径到叶子，维护树节点 \\
 IDGenRec & 学习的短文本ID & 交替更新ID生成器与推荐器 & 唯一性检查、前缀约束 \\
 TransRec & 数字ID、标题片段、属性 & 分方面生成监督 & 子串索引及多方面聚合 \\
 ETEGRec & 协同表示的残差量化码 & 编码与生成交替对齐 & 收敛后冻结，同步映射版本 \\
 RPG & 乘积量化的长代码 & 各代码位并行预测 & 合法ID图上的候选搜索 \\
 LIGER & 语义码与文本向量 & 代码预测及稠密匹配 & 生成集并入新物品，再稠密比较 \\
 COBRA & 稀疏码与稠密向量 & 级联代码及条件向量目标 & 分支内近邻，融合束与匹配分数 \\
 PLUM & 多模态与协同语义码 & 语言继续预训练、反馈加权微调 & 下一物品束搜索，代码解析 \\
 OneRec-V2 & 三级语义码与长上下文 & 上下文复用、真实反馈对齐 & 单物品代码解码及目录核验 \\
 \bottomrule
 \end{tabularx}
 \caption{标识生成模型的训练与候选接口。除RPG的各位并行预测外，这些实现主要使用逐步代码生成；混合方法还保留稠密表示与访问。这里的“联合”或“端到端”应按所连接的模块解释。}
 \label{tab:rs-id-models}
\end{table*}

% Outline: 3.2.2.1
\subsubsection{TIGER：语义标识生成}
\label{sec:rs-tiger}

任意数字ID没有表达商品之间的关系。\term{TIGER}先编码物品内容，再用残差量化得到\term{语义标识}（semantic ID, SID），使内容近似的对象能够共享部分代码\scite{rs-r06}残差量化的含义是每一级只编码前几级尚未解释的部分：令内容潜在向量为$\bm r_0$，第$\ell$级码本为$\{\bm e_{\ell,k}\}$，则
\begin{equation}
 c_\ell=\arg\min_k\|\bm r_{\ell-1}-\bm e_{\ell,k}\|^2,\qquad
 \bm r_\ell=\bm r_{\ell-1}-\bm e_{\ell,c_\ell}.
 \label{eq:rs-tiger-rq}
\end{equation}
编码器、码本和重建解码器以内容重建与量化相关损失训练，之后固定物品编码，让序列模型按式~\eqref{eq:rs-id-generation}预测下一物品代码。

作为二维例，内容向量$(0.9,0.2)$在第一级选择码向量$(1,0)$，剩余$(-0.1,0.2)$；第二级若恰有该向量，就能补齐表示。另一个背包可能选相同的首码，却在后级被区分。若两件商品最终代码均为$(2,5,1)$，原方法会增加消歧词元，成为$(2,5,1,0)$与$(2,5,1,1)$，并保存唯一映射；语义相近并不能代替身份区分。

请求读入历史物品的代码序列，束搜索产生下一代码，再通过映射取回真实物品。原方法可增大束并过滤无效代码；约束解码也是可选的合法性实现，不能把它自动当成所有TIGER设置的既有步骤。一个合法代码还可能指向下架商品，资格仍须检查。新物品若具有可用内容，可以编码，但能否在既定解码策略下被生成还取决于训练、消歧和候选覆盖。与目录检索比较时，应固定内容、历史与名额。

\Needspace{110mm}
% Outline: 3.2.2.2
\subsubsection{OneRec-Think：文字推理与物品生成}
\label{sec:rs-onerec-think}

当用户表达“这次想看轻松的内容”时，仅按过去ID续写很难直接处理这条文字条件。\term{OneRec-Think}连接物品词元与自然语言，在历史和文本条件下生成理由，再以理由作为生成下一物品标识的条件\scite{rs-r09}形式上可分解为$p(r,c\mid\bm x)=p(r\mid\bm x)p(c\mid\bm x,r)$，其中$r$是生成的文字路径，不是已知的偏好事实。

训练先通过物品—文本的多任务建立对应，再构造理由监督：选取与目标相关的历史形成较清晰的理由示范，训练模型从原始历史生成理由及目标物品。反馈优化针对推荐的稀疏命中问题，对每条理由计算
$R=\max_{\hat c\in B}\sum_{\ell=1}^m\mathbb I[\hat c_\ell=c_\ell^+]$，
其中$B$是该理由下的束候选。若目标码为$(2,5,1)$，束中最佳匹配为$(2,7,1)$，奖励便为2，即使完整物品尚未命中。组相对策略优化（GRPO）再比较同组理由的奖励、更新策略；这里的组相对优势是减去同组基准后的相对表现。匹配反馈更密集，但仍依赖训练目标物品，代码位相同也不证明对象满足文字要求。

假设历史以动作片为主，而当次文字要求轻松内容。理由可以归纳“保留冒险元素，降低紧张程度”，随后生成轻喜剧候选；检查应分别确认模型是否读到当前要求、候选是否符合描述、标识是否查得到。理由写得通顺不证明它忠实解释了用户为何喜欢。论文还提出提前计算理由与部分物品前缀、请求时完成解码的架构，须额外考虑历史陈旧和新要求能否及时进入。候选融合见第~\ref{sec:rs-candidate-fusion}~节，理由是否展示及其可信性见第~\ref{chap:rs-presentation}~章。

% Outline: 3.2.2.3
\subsubsection{DCGR：决策条件下的标识生成}
\label{sec:rs-dcgr}

同一生成器还可以先预测较短的结果类别，用它调节候选分布。OneTrans-V2中的\term{决策条件生成召回}（decision-conditioned generative retrieval, DCGR）在物品码前加入购买、消费金额、新类别探索与自然／赞助供给等结果维度\scite{rs-r43}这里的“决策”指模型预测的交互结果类别，不是用户明确选择，也不是自然语言理由。

原方法把购买、探索和赞助维度作为并行结果预测，组合其嵌入形成一个短词元；消费金额还依赖前面的购买信息，放在后续位置。令结果前缀为$z$，物品代码为$c$，训练最小化$-\log p_\theta(z,c\mid\bm x)$，标签从日志构造。服务预测前缀与物品，不读取未来真实结果。业务偏置$\Delta(z)$可加到前缀对数分数上，再与后续代码分数共同执行束搜索，改变目标类别的优先级。

例如两个假设前缀的概率为0.6与0.4，分别对应常规购买和新类别探索。给探索前缀增加$\log2$的偏置后，其未归一化权重从0.4变成0.8，高于前者0.6；后续候选仍要乘上各自条件物品概率。一个前缀得分高却没有高分物品时，联合束也可能舍弃它。偏置是软引导，代码碰撞、对象资格和类别覆盖仍须核验。完整共享架构见第~\ref{chap:rs-ranking}~章，赞助候选还需满足第~\ref{chap:rs-advertising}~章的竞价与预算条件。

\Needspace{85mm}
% Outline: 3.2.2.4
\subsubsection{LC-Rec：语言与协同语义对齐}
\label{sec:rs-lc-rec}

内容近似的两个商品可能被不同群体消费，语言模型也未必理解新增物品码。\term{LC-Rec}以文本表示的残差量化建立标识，在末级引入均匀语义映射，减少同组对象集中到少数代码的问题，再通过多类指令对齐语言与交互关系\scite{rs-r62}

主要任务仍是历史代码到下一代码。显式对齐让模型由标题、描述生成代码，也由代码恢复标题与描述；非对称预测则改变历史端或目标端的表示，例如历史代码预测目标标题、历史标题预测目标代码。基于评论提取的意图与历史归纳的偏好还可构造辅助任务，但它们是离线提取的监督，不应写成用户当次亲口提出的要求。各任务都可用目标词元负对数似然训练。

假设两款同为“防水背包”的商品，一款经常出现在通勤交互中，另一款主要与登山用品共同出现。内容编码可能给它们相近前缀，下一物品任务则让历史条件学会选择不同后缀；“代码到描述”任务维持语言对应。服务仍从真实历史生成标识、映射候选，并核验唯一性和资格。对齐成功需要分别检查代码检索与文字理解，不能由一个语言示例推断最终列表体验。

% Outline: 3.2.2.5
\subsubsection{LETTER：协同信号与标识分配}
\label{sec:rs-letter}

若量化器只为重建内容服务，生成器可能接收到不适合偏好区分的代码。\term{LETTER}把协同表示引入物品编码目标：内容重构保留语义，协同对齐使量化表示靠近已有推荐模型的表示，多样性正则改善代码嵌入分布与分配使用\scite{rs-r63}这些目标影响代码如何划分，随后再固定编码结果训练生成器。

设两款外观相近的背包具有相近内容向量，但协同向量分别接近通勤和登山群体。只有重构时，它们可能共用很多代码；加入协同对齐后，量化器有理由在后级区分它们。若多数物品挤在同一首码，代码利用率和训练样本分布也需一起检查。多样性约束针对代码空间，不等于承诺推荐列表一定覆盖多类需求。

生成阶段采用带温度的词元预测目标。温度降低时，高分竞争词元的作用更集中，形成原文讨论的排序引导训练；这仍是代码位上的竞争，不应解释为已经直接优化完整展示顺序。例如正确词元logit为2、对照为1，温度1时正确概率约0.731，温度0.5时约0.881，梯度尺度和困难对照权重也随之改变。服务通过代码解码和目录映射取得候选。额外协同模型的训练、量化器版本及映射刷新都应计入成本。

% Outline: 3.2.2.6
\subsubsection{SEATER：平衡树标识与层次学习}
\label{sec:rs-seater}

共享量化码本之外，还可让标识直接表示树路径。\term{SEATER}对协同物品向量递归执行有规模约束的聚类，形成平衡多叉树，为节点分配标识；历史由这些路径表示，生成器逐层预测下一物品路径\scite{rs-r64}树的同一层对应相近的分组粒度，合法后继由目录结构确定。

沿图~\ref{fig:rs-tree-retrieval}的八件物品教学二叉树，可以先分为户外与数码两组，再各分两组，最终到八个叶子。例如背包的路径为$(A,A_1,a)$，水壶为$(A,A_1,b)$，相机为$(B,B_1,e)$。预测首层$A$后，只能沿$A$的后继继续，束宽限制了同时探索的分支。与TDM逐层给节点评分相比，此处使用路径生成概率；两者都可能因早期裁剪遗漏叶子。

训练除目标路径似然外，还以层次对比学习约束节点表示，并用路径排序目标保持相似路径的次序。代码不是任意字符串，其关系应与建树和训练一致。内部节点和叶节点的标识、嵌入仍随物品数增长，不能套用少量共享量化码本的存储结论。新物品插入、删除或重新聚类会影响路径与历史编码，平衡树也需要维护，而不是天然适应动态目录。

% Outline: 3.2.2.7
\subsubsection{IDGenRec：可学习文本标识}
\label{sec:rs-idgenrec}

自然语言词表提供了已有语义，但商品标题常常太长或重名。\term{IDGenRec}让ID生成器从物品资料生成简短文本ID，并与推荐器交替训练，使ID既能描述对象，也有助于下一物品预测\scite{rs-r65}文本仍承担唯一寻址职责，不能把“读起来合理”当作唯一性成立。

推荐器更新时，先按当前ID生成器编码目录，将历史文本ID填入输入，以真实下一ID的词元似然训练。更新ID生成器时固定推荐器，但离散词元选择不能直接反传，于是用词元分布形成软嵌入：
\begin{equation}
 \widetilde{\bm e}_\ell
 =\sum_{v\in V}p_\phi(v\mid\text{物品资料},c_{<\ell})\,\bm e_v.
 \label{eq:rs-idgenrec-soft}
\end{equation}
推荐损失可沿该连续组合传回ID生成器。这是训练近似，服务仍生成离散文本，并检查其目录映射。

若两件商品都名为“城市背包”，文本ID可被分配为“轻量通勤”与“摄影通勤”。若仍重复，原方法增加候选生成的差异性或允许更长ID，直到得到未占用的标识。服务以包含全部有效ID的前缀树约束解码，避免生成无法解析的词串。每轮ID变化后，训练输入、前缀树和映射都要同步；跨域实验说明迁移的可能条件，不意味着文本生成器会自动维护新旧目录。

% Outline: 3.2.2.8
\subsubsection{TransRec：多方面标识与目录映射}
\label{sec:rs-transrec}

唯一ID擅长区分身份，标题和属性则提供可共享的语义。\term{TransRec}分别以数字ID、标题和属性组织历史及目标，训练模型按指定方面生成，再把多方面输出汇合到目录物品\scite{rs-r66}标题训练允许预测不同长度的子串，属性任务预测具体属性，数字ID任务保持身份对应。

服务用FM-index限制可生成的片段，使标题可以从相关子串开始，而不必从标题第一个词开始。这里的FM-index是压缩文本索引，与下一章的因子分解机FM无关。生成片段按其在物品标识中的覆盖关系落回对象；由于短串天然更容易获得高概率，原方法结合无条件出现频率修正片段分数，再做方面内与方面间聚合。

例如商品$a,b$都叫“城市背包”，但$a$属性为轻量，$b$属性为摄影隔层。标题片段会同时命中两件商品，属性“摄影隔层”只给$b$证据，数字ID又能直接指定$b$。假设三方面经各自校准后的教学分数对$a$为$(0.1,0.5,0)$，对$b$为$(0.3,0.5,0.4)$，在本例等权汇总规则下，$b$领先。这个算例只说明聚合接口，实际必须沿所用版本的频率与覆盖规则计算，不能随意相加原始生成概率。最终仍需去重和资格核验。

% Outline: 3.2.2.9
\subsubsection{ETEGRec：标识编码与生成的联合学习}
\label{sec:rs-etegrec}

量化器一旦固定，推荐损失就无法纠正不合适的代码分配。\term{ETEGRec}以协同物品表示为输入，连接残差量化器和生成推荐器，让编码随推荐目标调整\scite{rs-r67}它通过序列表示与目标物品在各级代码上的分布对齐，以及偏好表示与量化语义的对比学习，把两个模块的目标关联起来。

训练交替固定一侧、更新另一侧。量化器阶段不仅重建物品输入，还接收上述对齐约束；生成器阶段使用当前代码学习下一物品。量化器收敛后固定，再完成生成模型训练，减少代码持续变化造成的学习不稳。这里的端到端范围是标识编码和生成学习，尚未包含最终列表、曝光决策或广告支付。

例如背包的旧代码为$(2,5,1)$，更新后变为$(2,7,1)$，对应的历史输入、目标标签和代码—物品映射都必须切到同一版本。若训练已经学新码，服务仍查旧表，合法词元也可能解析失败。服务从历史生成当前版本代码再查目录。协同输入可以改善推荐区分，但并不自动赋予自然语言理解能力；若新增对象尚无协同表示，还需额外规定初始化与更新。

% Outline: 3.2.2.10
\subsubsection{RPG：长标识的并行预测}
\label{sec:rs-rpg}

增加单物品代码长度可以保留更多信息，却会增加逐词元解码次数。RPG用优化乘积量化把物品表示分到多个子空间，形成较长的代码；历史物品的多个代码嵌入先聚合，再由多个输出头同时预测下一物品各代码位\scite{rs-r68}其分解为
$p(c_1,\ldots,c_m\mid H)\approx\prod_{\ell=1}^m p_\ell(c_\ell\mid H)$，
训练求各位交叉熵之和，不在预测后位时读取前位的生成结果。

这相当于给定历史后各代码位条件独立的建模假设。若四位目标码的概率依次为0.8、0.6、0.7、0.5，联合得分为0.168；另一合法码的四项概率乘积为0.15，则前者领先。但把各位最大概率词元直接拼接，可能得到目录中根本不存在的组合。RPG在合法标识组成的相似图上搜索，以节点邻居的联合得分引导候选扩展，避免任意拼接落空。

服务的并行性属于一个物品的代码位，不等于一次共同决定多个展示位置。长代码会增加输出头与词表比较，图的构建、存储、搜索轮次也不能忽略。条件独立会限制代码位之间的直接依赖表达，实际收益应在相同候选数和目录协议下与逐步生成比较。

% Outline: 3.2.2.11
\Needspace{100mm}
\subsubsection{LIGER：生成候选与稠密匹配}
\label{sec:rs-liger}

新物品即使已有内容向量，也可能从未在代码生成训练中出现。\term{LIGER}在历史中同时使用语义码和文本表示，联合学习下一代码及目标文本向量匹配，再把生成与稠密分数接起来\scite{rs-r69}稠密目标用历史输出与目标文本向量的相似度构造softmax竞争，代码目标仍遵循式~\eqref{eq:rs-id-generation}。

服务先用束搜索恢复生成候选$C_{\mathrm{gen}}$，显式并入已知冷启动集合$I_{\mathrm{new}}$，再按稠密相似度比较两者的并集。假设束搜索得到$a,b,c$，新商品$d$没有被生成，但其内容匹配分数为0.9，高于$a,b,c$的0.8、0.7、0.6，则并集上的Top-2包含$d,a$。改善来自新物品真的进入比较，不能仅靠训练增加一个稠密损失就自动实现。

原方案假设待并入的冷启动集合相对较小。如果每天新增对象极多，直接全部加入会改变计算代价；还需另行设计新物品访问入口。固定文本向量仍要为目录保存，生成模型也没有消除这部分存储。该研究适合比较生成与稠密检索的收益和成本边界，不提供无条件的目录更新保证。

% Outline: 3.2.2.12
\subsubsection{COBRA：稀疏生成与稠密访问}
\label{sec:rs-cobra}

离散代码可以快速确定较粗的范围，范围内仍需要连续表示区分细节。\term{COBRA}在历史中组织稀疏标识与稠密表示，先预测稀疏代码，再以历史和该代码为条件预测稠密向量，最后在代码对应分支内检索近邻\scite{rs-r70}训练相应连接代码似然与条件稠密匹配，保留两种表示的不同职责。

例如束搜索保留分支$A,B$，条件向量在$A$内找到$a,b$，在$B$内找到$c,d$。如果简单只取最大概率分支，就看不到$c,d$。BeamFusion先对经过温度缩放的束分数作softmax，再对分支内经过缩放的余弦匹配分数作softmax，将两项相乘后统一比较。假设分支权重为0.6和0.4，$a$在$A$内的归一化匹配为0.5，$c$在$B$内为0.9，则融合分数分别为0.30与0.36，低权重分支的$c$仍能领先。分支内归一化常数不同，不能直接把原始余弦与束对数分数相加来替代此计算。

候选预算同时分配给束和分支内近邻，更多分支未必意味着更多不同需求，可能只是增加重复。代码映射、分支索引和稠密表示需保持版本一致。COBRA整合的是候选表示与访问；给候选打完整响应分数、组织展示以及处理广告资格和竞价仍有各自职责。

\Needspace{120mm}
% Outline: 3.2.2.13
\subsubsection{PLUM：语言预训练的生成召回适应}
\label{sec:rs-plum}

采用Transformer并不等于利用了语言预训练。PLUM从已有语言模型出发，扩展物品词元，以视频资料、物品标识、用户历史和普通文本继续预训练，再以反馈加权的目标标识预测适应工业召回\scite{rs-r72}其SID-v2结合多模态与协同共现，并通过多分辨率码本和渐进遮掩组织物品信息。

继续预训练让新增物品词元与已有语言上下文发生联系；召回微调则以具体用户历史预测后续物品，按行为反馈给样本不同权重。若两次目标观看的教学权重为1和0.2，则目标词元损失分别以1和0.2倍计入；权重改变的是训练贡献，不是宣称前者真实价值恰为后者五倍。历史还需保留观看状态，使完整观看与快速跳过不被编码成相同事件。

服务读取当前历史，通过束搜索输出多个下一物品标识，再映射回真实视频。候选之间并没有因此自动建立共同列表目标。语言初始化带来的收益应与领域继续预训练分开消融，成本也包括这两个阶段；代码碰撞、无效生成和目录持续变化仍需处理，不能把模型能说自然语言当作它已满足个性化召回。

\Needspace{145mm}
% Outline: 3.2.2.14
\subsubsection{OneRec-V2：上下文复用与反馈对齐}
\label{sec:rs-onerec-v2}

长历史的处理与短物品码的解码具有不同计算形状。\term{OneRec-V2}让上下文处理器把资料和历史变为可共享的键值表示，短解码序列通过交叉注意读取条件，再逐步预测目标物品的三级标识。其主体仍是单物品标识预测，不能沿用OneRec早期版本的整会话输出假定\scite{rs-r73}

在反馈对齐阶段，原报告使用相似时长视频中的用户播放分位数，并结合负反馈构造奖励。分位数缓解直接比较不同视频长度的尺度问题，但仍是短期行为代理。梯度有界策略优化（gradient-bounded policy optimization, GBPO）调整概率比值的分母与梯度，使更新受控；它不同于仅拟合一个固定奖励模型后无限放大高奖励输出。

具体地，记当前物品生成概率为$p_\theta$，日志旧概率为$p_{\mathrm{old}}$，优势为$A$。原目标使用$-A p_\theta/d$，正优势时$d=\max(p_{\mathrm{old}},\operatorname{sg}(p_\theta))$，负优势时$d=\max(p_{\mathrm{old}},1-\operatorname{sg}(p_\theta))$；$\operatorname{sg}$表示分母中的当前概率停止梯度。若$p_{\mathrm{old}}=0.02,p_\theta=0.04,A=1$，比值从普通重要性比的2变成1，但由于分母停止梯度，样本仍提供提高概率的梯度；若$A=-1$，分母为0.96，比值约0.0417，负例的过大梯度受到限制。

例如某次真实生成曝光的日志保存了物品标识及其生成概率，之后观察到较高播放分位数。训练可以把这次结果与原策略概率对应，调整该标识的生成概率；若改用传统多阶段流程的曝光，通常没有与该生成模型一致的真实日志概率。原文在这部分使用当前概率作为代理，因而不能把混合训练目标解释成一般的无偏离策略评价。上下文缓存还须及时更新，播放奖励也不能直接证明长期留存或用户福祉改善。

% Outline: 3.3
\section{候选融合与维护}
\label{sec:rs-candidate-management}

单个模型输出高分对象，还没有形成可交给后续阶段的候选。不同路径会重复命中，某件物品可能刚刚缺货，某个新物品可能还没进入索引。本节把集合的覆盖分配与对象的状态核验分开：前者决定有哪些比较机会，后者确认这些机会在当前请求中有效。两者可以在访问前后交错执行，不要求固定为一次线性处理。

% Outline: 3.3.1
\subsection{多路融合与覆盖}
\label{sec:rs-candidate-fusion}

设路径$r$产生$C_t^{(r)}$，一种基础融合是先取并集，再按名额或统一分数截断。路径可以来自热门先验、历史邻域、内容、协同、当前查询或探索，不应仅按算法名称分配固定份额。判断一条路径是否有用，要看它增加了哪些其他路径没有提供的有效对象，以及这些对象满足什么需求。

假设历史路径返回$\{a,b,c\}$，内容路径返回$\{b,c,d\}$，新供给路径返回$\{d,e\}$。原始结果有8条，去重后只有5件。前两路已覆盖$\{a,b,c,d\}$，第三路的额外覆盖仅为$e$；重复命中$d$可作为来源证据，但不能再次记为新增物品。如果总名额为4，必须决定是否给$e$保留机会，以及谁被挤出。简单按返回顺序截断，会把路径顺序变成隐含的配额政策。

多兴趣方法还需要避免一个兴趣的近邻填满所有名额。可以为兴趣保留最小配额，再用剩余名额比较分数；也可以像ComiRec那样在候选上计算相关性与类别差异的边际收益。两者都要区分“类别标签不同”与“覆盖不同需求”。同为背包的通勤款和登山款可能覆盖两种任务，类别不同的相机与镜头也可能只服务同一种任务。

新供给和长尾物品可以依赖内容入口、先验或探索配额进入比较，但候选机会还不是测试曝光。若它们每次进入$C_t$后都被下一阶段截掉，就没有产生新的真实反馈。第~\ref{chap:rs-presentation}~章将讨论展示探索；此处应记录路径来源、去重前后名额、有效新增候选及后续存活率，才能区分“路径没有召回”与“下游没有选择”。

候选融合还改变了排序模型所面对的分布。增加一条语义相近但行为稀少的路径，会让精排遇到训练中罕见的物品；观察到最终指标变化时，不能只归因于召回命中率。路径消融应固定总名额与资格条件，并同时检查兴趣、长尾覆盖及计算成本。增加$K$可以机械地增加命中，只有在同等预算下仍有更有用的覆盖，才说明融合更适合当前任务。

% Outline: 3.3.2
\subsection{候选核验与目录维护}
\label{sec:rs-candidate-maintenance}

候选核验检查对象是否真实存在、属性是否满足请求、库存或供给是否可用、时间范围是否有效。“防水通勤背包，价格不超过500元”至少要求正确解释属性与价格限制；模型高分不能覆盖不满足条件的事实。广告还要补充地域、人群、投放状态等资格，搜索则要检查显式查询约束，分别见后续专题章。

资格可以先过滤，也可以在检索后核验。先过滤能减少无效比较，但索引未必支持全部复杂条件；后过滤更灵活，却可能使原本100个结果只剩20个。这时需要续取、更换访问范围或允许返回较少有效候选。评价分母也必须固定：把无资格对象当成应召回目标，会错误惩罚合法模型。

目录维护负责让这些依据跟随对象变化。新物品需要资格记录、内容或协同表示、索引条目或生成代码映射；旧物品失效时，即使向量或代码仍在，也应立即通过资格层阻止错误输出。属性变化可能要求重编码，而库存短时变化通常更适合单独维护，未必每次重训表示。对于树和生成码，重新分组还会改变路径或代码语义，模型与解析表必须使用兼容版本。

例如新背包$e$在10:00加入目录，资格库10:01可见，向量索引10:03更新，标识映射10:05更新，生成器又未必能马上给该码足够概率。不同延迟分别导致“不能访问”“能解析但很难生成”等现象。失效商品也可能在索引中短暂残留，只要当前资格核验及时，它就不该进入最终候选。维护评价应报告新增后的遗漏率、失效候选率和恢复时间，而不仅是每小时完成了多少次更新；缓存和索引服务的具体实现归系统卷。

% Outline: 3.4
\section{召回学习与评价}
\label{sec:rs-retrieval-learning}

候选访问方式确定后，仍需知道模型从哪里学到“值得比较”。已交互物品通常可作为训练正例，未观测物品只表示日志没有该条记录。它可能没有被展示，也可能被展示但未响应，甚至是用户尚未发现的合适对象。把这些状态全部称为“不喜欢”会使训练目标的解释超出证据。

% Outline: 3.4.1
\subsection{召回训练与表示学习}
\label{sec:rs-retrieval-training}

大目录目标常需要近似竞争集合。随机采样用少量对象降低输出成本，\term{困难负例}把计算集中在当前容易混淆的对象，批内负例复用其他样本的正物品。它们改变的是训练比较分布，不是自动获得真实拒绝标签。标识生成主要在代码词表上预测，向量恢复可以不使用推荐负例，两种接口不必强行共享同一采样程序。

辅助表示学习提供另一类信号。遮蔽已有物品、匹配属性或比较同一历史的不同视图，可以从现有记录构造额外任务；用“下一目标相同”选择正例则已经使用推荐标签。辅助任务的成功不能证明曝光偏差已消除，但它们可以在相同数据下改变编码器学到的结构。

% Outline: 3.4.1.1
\Needspace{125mm}
\subsubsection{采样softmax与物品频率修正}
\label{sec:rs-sampled-softmax}

式~\eqref{eq:rs-directory-softmax}的分母包含全目录。若每次只比较采样物品，而且热门物品被抽到的概率更高，未经修正的训练就会让它们承担更多负例竞争。YouTube的采样训练以及后续频率校正工作研究了如何把这种非均匀取样与目标目录分布联系起来\scite{rs-r01,rs-r04}

先看一个便于推导的估计器。固定正例$i^+$，从其余目录按$q(j)>0$独立有放回抽取$m$个负例$J_1,\ldots,J_m$，则归一化常数可估计为
\begin{equation}
 \widehat Z=\exp s_{i^+}
       +\frac1m\sum_{k=1}^{m}\frac{\exp s_{J_k}}{q(J_k)}.
 \label{eq:rs-softmax-estimator}
\end{equation}
对抽样取期望后，负例部分恢复为$\sum_{j\ne i^+}\exp s_j$，所以$\widehat Z$是该常数的无偏估计。但$\log\widehat Z$和由比值形成的梯度并不因此在有限样本下无偏。分数修正$s_j-\log(mq(j))$只是把对应的重要性权重移入指数；正例是否修正、重复是否去重、采样是否有放回，都属于具体估计器的约定。

批内负例来自其他样本目标，频率与数据流相关。该频率校正工作主要处理流式批内对象，以相邻出现的训练步数间隔估计物品进入随机批次的概率$p_j$，对包括正例在内的各竞争logit使用$s_j-\log p_j$；哈希记录与滑动更新使估计能适应变化。若另外混入随机采样对象，需要按实际混合来源重新确定修正量，不能直接套用批内频率。作为简化示例，两物品原分数均为1，采样概率分别为0.8与0.2，修正分数为1.223与2.609。长尾对象每次被采到时权重更大，用来补偿它出现次数较少，不代表应该在部署中直接给所有长尾物品加同样奖励。

训练完后，服务通常按模型原分数访问目录，不能把训练用的$-\log q$不加区分地叠加到线上分数。采样频率修正也不同于系统曝光去偏：它知道的是训练器怎样抽对照，不一定知道历史策略为何展示物品。评价时只抽少数候选又是第三件事，会改变测试排名难度，见第~\ref{sec:rs-retrieval-evaluation}~节。

\Needspace{85mm}
% Outline: 3.4.1.2
\subsubsection{ANCE全目录困难负例训练}
\label{sec:rs-ance}

随机负例常与正例差得很远，模型很快就能区分。\term{ANCE}让当前编码器的近邻索引提供困难负例：先编码语料并建索引，训练时按查询取回高相似对象，用已知正例与这些对象竞争，再异步刷新目录向量和索引\scite{rs-s03}训练器因而能看到比同一小批次更广的混淆对象。

原任务是稠密文本检索。迁移到推荐时，必须另定正交互、物品内容或ID表示，以及哪些对象在该请求下有资格。假设正例为通勤背包，随机负例为电饭煲，当前分数分别为0.8和0.1；索引找到登山包，分数0.75，与正例更接近，因而提供更强区分压力。但是登山包可能同样有用，未标注不等于不相关，困难负例中也可能含误负例。

异步索引降低同步重编码的代价，却会使用稍旧的表示；刷新间隔改变负例难度和分布。训练应记录索引版本、已知正例排除规则及物品频率，避免把旧索引或热门偏移误解为方法本身的收益。ANCE检索的是训练对照，前章历史检索选的是已发生行为，本章线上召回取的是待比较物品，三者的访问目标不同。

% Outline: 3.4.1.3
\subsubsection{S3-Rec：多关联的序列预训练}
\label{sec:rs-s3rec}

下一物品标签稀疏时，已有序列和属性仍能提供更多关联。\term{S$^3$-Rec}在正式推荐训练之前学习四类任务：物品与其属性的关联预测、上下文恢复遮蔽物品、上下文预测被遮蔽物品的属性，以及上下文恢复一段行为片段。各任务通过区分正确关联与对照关联，使编码器保留属性、物品、片段和序列之间的信息\scite{rs-r81}

在“登山包—水壶—相机”的教学序列中，登山包与“户外”属性提供物品—属性关联；遮住水壶后，其余上下文用于识别真实水壶及它的属性；遮住“登山包—水壶”片段，则要求上下文与正确片段的表示匹配。四者的预测对象不同，不能把它们都缩写成“掩码物品预测”。目标通常联合加权预训练，再以真实后继监督微调下一物品评分，得到的向量或目录分数才用于召回。

预训练可以在已经观察完整的历史上使用双向上下文；后续下一物品任务必须恢复因果约束。对时点$t$作评估时，预训练数据本身也不能偷偷包含$t$之后的交互。辅助任务只重用了已有资料，没有制造未发生的用户反馈；属性缺失或错误还可能使预训练约束偏离需求。

\Needspace{75mm}
% Outline: 3.4.1.4
\subsubsection{CL4SRec：增强视图的序列对比学习}
\label{sec:rs-cl4srec}

同一段历史不应因偶然噪声而产生完全不同的表示。\term{CL4SRec}通过裁剪、遮蔽或局部重排生成两个视图，让共享编码器给出接近的表示，并与下一物品损失联合训练\scite{rs-r82}对锚点$\bm h$、正视图$\bm h^+$及对照集合$B$，一个对比方向为
\begin{equation}
 \ell_{\mathrm{con}}
 =-\log\frac{\exp(\operatorname{sim}(\bm h,\bm h^+)/\tau)}
 {\exp(\operatorname{sim}(\bm h,\bm h^+)/\tau)
    +\sum_{\bm h^-\in B}\exp(\operatorname{sim}(\bm h,\bm h^-)/\tau)}.
 \label{eq:rs-sequence-contrast}
\end{equation}
温度$\tau>0$控制相似度区分尺度，另一视图也可作为锚点计算对称目标。训练后的编码器沿既定下一物品评分接口召回，服务不需要反复增强请求历史。

若一次相机点击是偶然浏览，遮掉它可能仍保留“通勤装备”需求；若遮掉唯一一次通勤包购买，却留下户外浏览，视图就可能改变主要意图。裁剪过短、重排改变因果顺序、同批用户实际具有相似偏好，都会削弱正负视图假设。对比损失下降只说明模型更满足所选不变性，是否改善真实候选覆盖还需独立评价。

% Outline: 3.4.1.5
\subsubsection{DuoRec：模型与同目标视图的对比学习}
\label{sec:rs-duorec}

直接改动序列可能删去关键意图。\term{DuoRec}保留物品顺序，用不同dropout得到同一输入的两种编码结果；另外从训练集中选择下一目标相同的历史作为正例，以推荐损失和对比正则联合学习\scite{rs-r83}前一类是模型扰动视图，后一类依赖下一物品标签，不能笼统称为没有监督的信息。

例如用户甲的历史为“地铁通勤—轻量水壶”，用户乙的历史为“雨衣—电脑内胆包”，二者下一条训练交互都指向同一防水背包。模型可以把两段历史的表示拉近，帮助区分这个共同目标与其他物品；训练时再配合各自序列的dropout视图。服务只编码请求时已有历史，再按物品评分召回，不需要知道真实下一物品，也不能利用测试目标去挑正例。

同目标不表示两个人所有偏好相同，高频物品会形成很大的正例组，冷门物品则可能没有可配对历史。因此应检查表示是否仍能区分不同目标、不同频率分组是否受益，以及对比正例是否在时间切分内合法。DuoRec针对表示集中和区分能力提出训练约束，并没有证明同目标用户可以在所有任务上互换。

% Outline: 3.4.2
\subsection{召回效果评价}
\label{sec:rs-retrieval-evaluation}

召回评价必须先固定可用目录$I_t$、请求信息、候选名额$K$和计算预算。设可用于评价的目标集合为$G_t\subseteq I_t$，则候选召回率为$|C_t\cap G_t|/|G_t|$；单目标时对应命中指示。若$G_t$来自留出的未来交互，这个指标只衡量对这些已观测目标的覆盖，不等于所有真实相关物品的召回率。没有目标的请求应另定统计规则，而不是默默从分母删除。

全目录与采样候选协议可以改变结论。假设目标在全目录排第1000，但测试只抽99件容易区分的对照，它可能在这个小集合中排第一。这样的命中不能证明模型能从海量目录召回目标。向量近邻还应区分模型评分目标与索引近似误差：相对于精确Top-$K$的索引命中率，和相对于用户目标的召回率，是不同指标。

多路方案需分别评价单路、去重融合后以及最终交给排序的集合。除目标命中外，还要报告兴趣或用途覆盖、长尾与新供给覆盖、无效对象比例、重复率、时延与资源成本。目录更新后再检查新对象出现延迟、失效对象残留与恢复时间。若一个模型依赖更多内容、更长历史或更多候选，比较时应说明这些条件，而非只列一个最终百分比。

候选被下游截掉时，可以追踪其分数和来源，但没有展示就通常没有对应的真实响应。不能把“被召回但未点击”一律标成负反馈，因为其中很多对象根本没有展示。日志的策略覆盖、实际曝光、随机化试验与离策略识别由第~\ref{chap:rs-evaluation}~章完整建立。历史行为检索则评价选出了哪些已有证据，不能与本章目录候选的命中混用。

% Outline: 3.5
\section{本章小结}
\label{sec:rs-retrieval-summary}

候选召回把需求依据转化为有限的比较空间。目录检索可以使用邻域、潜在因子、独立编码、多兴趣、图传播、树访问或生成的连续表示；标识生成则让模型预测能恢复真实对象的代码。多行为与社交关系改变证据，变分或扩散目标改变学习方式，只有最终怎样取得物品才决定候选接口。

多路融合补充覆盖，资格核验排除当前无效对象，目录维护让索引与映射跟上变化。训练对照和辅助任务并不等于真实拒绝，命中率也依赖目录、时间与测试协议。这些条件共同决定下一阶段究竟能看到哪些对象。候选有了比较机会后，还要估计给定曝光下的响应、学习偏好次序，并区分高响应与真正带来额外收益的机会；下一章从这些不同目标展开排序。
